\paragraph{3.2. Non-commutative root-influence inequality.}
For clarity in the proofs below, we remind our convention that all tuples or blocks are denoted with boldface fonts (e.g. $\mathbf U_1$ or $\mathbf A$), while a single element is denoted without boldface (e.g. $U_1(i)$ or $A_i$ or $A$). Before proceeding with the proof, we restate the statement for convenience.
\paragraph{Theorem 4 (Non-commutative root-influence inequality).}
Let $f$ be a homogeneous degree-$d$ block-multilinear form. Then, for blocks $b\in\{1,d\}$,
\[\|f\|_{\mathrm{cb}}\geq\frac1{\sqrt{e(d+1)}}\sum_{i=1}^n\sqrt{\operatorname{Inf}_{b,i}(f)}.\]
\begin{proof}[Proof of Theorem 4]
Since $f$ is homogeneous, we can write
\begin{align*}
f(\mathbf x_1,\ldots,\mathbf x_d)
&=\sum_{i_1,\ldots,i_d\in[n]}\widehat f_{i_1,\ldots,i_d}\,x_1(i_1)x_2(i_2)\cdots x_d(i_d)\\
&=\sum_{i=1}^nx_1(i)\underbrace{\left(\sum_{i_2,\ldots,i_d\in[n]}\widehat f_{i,\ldots,i_d}\,x_2(i_2)\cdots x_d(i_d)\right)}_{:=f_i(\mathbf x_2,\ldots,\mathbf x_d)}.
\end{align*}
In this case, it follows from (7) that for each $i\in[n]$, we have
\[\operatorname{Var}[f_i]=\|f_i\|_2^2=\operatorname{Inf}_{1,i}(f)\quad\text{and}\quad\operatorname{Var}[f]=\sum_{i=1}^n\operatorname{Inf}_{1,i}(f).\tag{8}\]
Let us denote the corresponding non-commutative block-multilinear polynomials by $f(\mathbf U_1,\ldots,\mathbf U_d)$ and $f_i(\mathbf U_2,\ldots,\mathbf U_d)$ where $\mathbf U_b=(U_b(1),\ldots,U_b(n))$ denotes the $b$th block of non-commutative variables. To show a lower bound on $\|f\|_{\mathrm{cb}}$ it suffices to exhibit a collection of square matrices $\{U_b(i)\}_{b\in[d],i\in[n]}$ with operator norm at most 1, such that $\|f(\mathbf U_1,\ldots,\mathbf U_d)\|_{\mathrm{op}}$ is large.
Applying Lemma 12 for the homogeneous case (with $p=f$, $q_i=f_i$ for $i\in[n]$, and $q_0=0$), it follows that for every $\delta>0$ there exists an integer $N$ and a choice of tuples $\mathbf U_1,\ldots,\mathbf U_d$ of $N\times N$ unitaries such that
\begin{align*}
\|f\|_{\mathrm{cb}}&\geq\|f(\mathbf U_1,\ldots,\mathbf U_d)\|_{\mathrm{op}}\\
&\geq\frac1{\sqrt{e(d+1)}}\sum_{i\in[n]}\|f_i\|_2-\delta\\
&\overset{(8)}=\frac1{\sqrt{e(d+1)}}\left(\sum_{i=1}^n\sqrt{\operatorname{Inf}_{1,i}(f)}\right)-\delta.
\end{align*}
Taking $\delta\to0$, we get the statement of the lemma. The proof for the inequality when $b=d$ is the last block follows similarly by using the right polar decomposition analogue of Lemma 12.
\end{proof}
\paragraph{3.3. Aaronson--Ambainis Conjecture for non-homogeneous forms.}
In this section, we prove Theorem 3, which requires handling non-homogeneous forms. The proof will be similar to the proof of Theorem 4 but we will need to be careful about certain details.
\begin{proof}[Proof of Theorem 3]
Any block-multilinear polynomial $f(\mathbf x_1,\ldots,\mathbf x_d)$ can be written as
\[f(\mathbf x_1,\ldots,\mathbf x_d)=\mathbb Ef+\sum_{b\in[d]}f_b(\mathbf x_b,\mathbf x_{b+1},\ldots,\mathbf x_d),\]
where $f_b$ consists of all monomials of $f$ that start with a variable in the $b$th block $\mathbf x_b$. Note that $f_b$ depends only on the variables in blocks $\mathbf x_b,\mathbf x_{b+1},\ldots,\mathbf x_d$. Moreover, it follows from Parseval that
\[\operatorname{Var}[f]=\sum_{b\in[d]}\|f_b\|_2^2=\sum_{b\in[d]}\operatorname{Var}[f_b],\tag{9}\]
so there exists a block $\beta\in[d]$ such that $\operatorname{Var}[f_\beta]\geq\frac1d\operatorname{Var}[f]$.
Since $f_\beta$ contributes a lot to the variance, it is natural to try to find an influential variable in the block $\mathbf x_\beta$. Towards this end, we pull out the variables $x_\beta(i)$ and write
\[f_\beta(\mathbf x_\beta,\ldots,\mathbf x_d)=\sum_{i\in[n]}x_\beta(i)f_{\beta,i}(\mathbf x_{\beta+1},\ldots,\mathbf x_d),\]
for block-multilinear polynomials $f_{\beta,i}(\mathbf x_{\beta+1},\ldots,\mathbf x_d)$. Note that some of the $f_{\beta,i}$'s could be identically zero, so let us define $S$ to be the set of those $i$ such that $f_{\beta,i}$ is non-zero. We note that
\[\|f_{\beta,i}\|_2^2=\operatorname{Inf}_{\beta,i}(f_\beta)\leq\operatorname{Inf}_{\beta,i}(f),\tag{10}\]
which implies (using Parseval's identity) that
\[\frac1d\operatorname{Var}[f]\leq\operatorname{Var}[f_\beta]=\sum_{i\in S}\|f_{\beta,i}\|_2^2=\sum_{i\in S}\operatorname{Inf}_{\beta,i}(f_\beta).\tag{11}\]
Denote the corresponding non-commutative block-multilinear polynomials by $f(\mathbf U_1,\ldots,\mathbf U_d)$, $f_b(\mathbf U_b,\ldots,\mathbf U_d)$, and $f_\beta(\mathbf U_{\beta+1},\ldots,\mathbf U_d)$ where $\mathbf U_b=(U_b(1),\ldots,U_b(n))$ denotes the $b$th block of non-commutative variables. To show a lower bound on $\|f\|_{\mathrm{cb}}$ it suffices to exhibit a collection of square matrices $\{U_b(i)\}_{b\in[d],i\in[n]}$ with operator norm at most 1 such that $\|f(\mathbf U_1,\ldots,\mathbf U_d)\|_{\mathrm{op}}$ is large.
We set the matrices in blocks $\mathbf U_1,\ldots,\mathbf U_{\beta-1}$ to be zero (that is, the all-zero matrix $0$). Note that with this choice all polynomials $f_b(\mathbf U_b,\ldots,\mathbf U_d)$ where $b<\beta$ vanish and the non-commutative polynomial becomes
\begin{align*}
f(0,\ldots,0,\mathbf U_\beta,\mathbf U_{\beta+1},\ldots,\mathbf U_d)
&=\sum_{i\in S}U_\beta(i)f_{\beta,i}(\mathbf U_{\beta+1},\ldots,\mathbf U_d)\\
&\quad+\sum_{b=\beta+1}^df_b(\mathbf U_b,\mathbf U_{b+1},\ldots,\mathbf U_d)+\mathbb Ef,
\end{align*}
which is a non-commutative polynomial of the form considered in Lemma 12 (with $m=|S|$, $q_i=f_{\beta,i}$ and $q_0=\sum_{b=\beta+1}^df_b+\mathbb Ef$). Thus, by Lemma 12 for every small $\delta>0$ there exists an integer $N$ and a choice of $N\times N$ matrices for the blocks $\mathbf U_\beta,\ldots,\mathbf U_d$ such that
\begin{align*}
\|f\|_{\mathrm{cb}}&\geq\|f(0,\ldots,0,\mathbf U_\beta,\mathbf U_{\beta+1},\ldots,\mathbf U_d)\|_{\mathrm{op}}\\
&\geq\frac1{\sqrt e(d+1)}\sum_{i\in S}\|f_{\beta,i}\|_2-\delta\\
&\overset{(10)}=\frac1{\sqrt e(d+1)}\left(\sum_{i\in S}\sqrt{\operatorname{Inf}_{\beta,i}(f_\beta)}\right)-\delta\\
&\overset{(11)}\geq\frac1{\sqrt e(d+1)}\left(\frac{\sum_{i\in S}\operatorname{Inf}_{\beta,i}(f_\beta)}{\sqrt{\operatorname{MaxInf}(f)}}\right)-\delta\\
&\overset{(10)}\geq\frac1{\sqrt e(d+1)^2}\left(\frac{\operatorname{Var}[f]}{\sqrt{\operatorname{MaxInf}(f)}}\right)-\delta.
\end{align*}
Taking $\delta\to0$ and using the assumption that $\|f\|_{\mathrm{cb}}\leq1$, we obtain the statement of the theorem:
\[1\geq\|f\|_{\mathrm{cb}}\geq\frac1{\sqrt e(d+1)^2}\cdot\frac{\operatorname{Var}[f]}{\sqrt{\operatorname{MaxInf}(f)}}\quad\Longrightarrow\quad\operatorname{MaxInf}(f)\geq\frac{(\operatorname{Var}[f])^2}{e(d+1)^4}.\]
\end{proof}
\paragraph{3.4. Approximating completely bounded forms with decision trees.}
In this section, we briefly mention how to obtain Corollary 5. Aaronson and Ambainis [1, Theorem 3.3] showed that querying the most influential variable reduces the variance of the function $f$, and if that influence is lower bounded by a polynomial in $\operatorname{Var}[f]/d$, then after $\operatorname{poly}(d)$ queries (the exact quantitative dependence can be read off from their proof), the variance of the function becomes small enough so that it can be approximated almost-everywhere by its expectation. Since the family of degree-$d$ block-multilinear forms with completely bounded norm at most one is closed under restrictions, one can apply Theorem 3 repeatedly. This gives us Corollary 5.
